\subsection{Семинары «Структурные паттерны» (6 очных часа)}

Обеспечьте то, что в вашей программе можно будет добавлять, удалять и изменять сущности (выбирая соответствующие пункты меню;
выбирать изменяемую или удаляемую сущность можно любым простым способом, например, по порядковому 
номеру при выводе на экран).

\begin{itemize}
   \item Используя паттерны <<Адаптер>> и <<Мост>> реализуйте два варианта консольного интерфейса: тот, что был раньше и
   цветной (библиотека rich). Выбор осуществляется передачей параметра командной строки (\verb|--colored |-- значит новый вариант);
   \item Реализуйте возможность объединять сущности, что уже были раньше, в папки (папки могут быть внутри папок). Для этого
   используется паттерн <<Компоновщик>>. Достаточно обеспечить возможность создавать новые папки и при создании сущности или 
   папки выбирать ту, в которой она будет находиться; 
   \item Реализуйте возможность протоколировать операции, что происходят внутри хранилища по выбору 
   пользователя (для этого используется паттерн <<Декоратор>>). Реализуйте две опции \verb|--log |\textit{имя файла} и 
   \verb|--log-json |\textit{имя файла}, которые можно использовать независимо, и осуществляющие протоколирование в обычном
   текстовом формате и в формате json соответственно
   \item Создайте класс, инкапсулирующий всю работу вашей системы, кроме ввода-вывода, чтобы можно было любые необходимые действия
   осуществлять посредством вызова методов этого класса (паттерн <<фасад>>)
   \item Дайте возможность указывать, что сущность в вашей программе (или папка), не доступны для изменения или удаления; контролируйте
   это с помощью паттерна <<заместитель>>
   \item Добавьте возможность выбирать цвет фона для отображения каждой сущности и папки, создаваемой пока цвет фона актуален.
   Актуальность меняется, когда пользователь выбирает новый цвет фона. Здесь используется паттерн <<Приспособленец>>
\end{itemize}
